\lecture{26}{Mon 14 Nov 2022 10:38}{}

\begin{fact}
	Every finite field has $p^n$ elements for prime $p$ and natural number $n$ ; and for every $p,n$ we have a finite field with $p^n$ elements. Furthermore, every field with $p^n$ elements are isomorphic.
\end{fact}

\subsection{Polynomial Rings}

Given $R$ a ring, we can define $R[x] = \{a_0 + a_1x + \ldots + a_n x^n: a_i \in R, n \ge 0\} $. We concentrate on situation when $R$ is a field.
\begin{enumerate}
	\item $F[x]$ forms a commutative ring with unit.
		 \begin{proof}
			 Consider $p(x) = a_0 + a_1x + .. + a_n x^n \in F[x]$, and $q(x) = b_0 + b_1x + \ldots + b_m x^m \in F[x]$, $n \ge m \ge 0$.

			 \begin{itemize}
			 	\item (equality)
			 $p(x) = q(x)$ iff $a_0 = b_0, \ldots, a_m = b_m$, $a_j = 0$ for $j > m$. 
		                \item (addition)
					\[
						p(x)+ q(x) = (a_0+b_0) + (a_1+b_1)x + \ldots + (a_m + b_m)x^m + a_{m+1}x^{m+1} + \ldots + a_n x^n
					.\] 
				\item (multiplication)
					Define $p(x) q(x) = c_0 + c_1x + \ldots + c_r x^r \in F[x]$ where  $c_i = \sum_{l=0}^{i} a^l b^{i - l}$, adopt convention that $a_i = 0$ for $i > n$, $b_j = 0$ for $j > m$.
				\item (Additive Identity) $0 \in F[x]$
				\item (Additive Inverse) $-(a_0+a_1x+\ldots+a_nx^n) = (-a_0) + (-a_1) x + \ldots + (-a_n)x^n$
				\item Assiciativity and Commutativity of addition inherited from $F$.
				\item (Multiplicative Identity) $1 \in F[x]$ since $1 \in F$
				\item (Commutativity of Multiplication) inherited from $F$.

			 \end{itemize}
		\end{proof}
\end{enumerate}

We measure the size of polynomials with ``degree'', which is analogous to $\log|n|$ for $n \in \mathbb{Z}$.

\begin{definition}[Degree of Polynomial]
	Suppose $f(x) = a_0+a_1x + \ldots + a_n ^n$ with $a_n \neq 0$, then the \textit{degree} of $f(x)$ is $\operatorname{deg} f = n$. By convention, $\operatorname{deg} (0) = -\infty $.
\end{definition}
\begin{example}
	One has $3x^2 + 2 \in \mathbb{Z}_5[x]$, and $\operatorname{deg} (3x^2 + 2) = 2$, and $\operatorname{deg} (2) = 0$. 
\end{example}
\begin{proposition}
	The following statements hold for $p(x), q(x) \in F[x]$ :
	\begin{itemize}
		\item $\operatorname{deg} \left( p(x) + q(x) \right) \le \max \{\operatorname{deg} p, \operatorname{deg} q\} $
		\item $\operatorname{deg} (pq) = \operatorname{deg} p + \operatorname{deg}  q$
	\end{itemize}
\end{proposition}
\begin{proof}
	exercise.
\end{proof}

\begin{lemma}
	When $F$ is a filed, then $F[x]$ is an integral domain.
\end{lemma}
\begin{proof}
	Need to check $F[x]$, which is commutative ring with unity, has no zero divisors. But if  $p(x) , q(x) \in F[x]$ and $p(x)q(x) = 0$, then
	\[
		- \infty = \operatorname{deg} \left( p(x)q(x) \right) =\operatorname{deg} p + \operatorname{deg} q
	.\] 
	Which is possible when either $\operatorname{deg} p = -\infty $ or $\operatorname{deg} q = -\infty $, or both. So one of them is $0$.
\end{proof}

\begin{theorem}[Division Algorithm]
	Suppose that $f(x), g(x) \in F[x]$ and that $g(x) \neq 0$. Then
	\[
		f(x) = q(x)g(x) + r(x)
	\] 
	where the \textit{quotient} $q(x)$ and \textit{remainer} $r(x) \in F[x]$, and $r(x) = 0$ or  $\operatorname{deg} r < \operatorname{deg} g$.
\end{theorem}
\begin{proof}
	If $\operatorname{deg} f < \operatorname{deg} g$, take $q = 0$ and $r = f$ and we are done. So we assume $\operatorname{deg} f\ge \operatorname{deg} g$, and say
	\begin{align*}
		f(x)&= a_0+a_1x + \ldots+a_mx^m \\
		g(x)&= b_0+b_1x+\ldots+b_nx^n & m\ge n\ge 0, a_m \neq 0, b_n\neq 0\\
	.\end{align*}
	Proceed induction and assume conclusion holds for $\operatorname{deg} (f) < m$ for some $m\ge 1$. Now assume $\operatorname{deg} f = m$ as above, and observe
	\begin{align*}
		&f(x) - \frac{a^m}{b^n}x^{m-n}g(x) \\
		&= \left( a_{m-1} - \frac{a_m}{b_n}b_{n-1 } \right) x^{m-1} + \ldots+ \left( a_{m-n}- \frac{a^m}{b^n}b_0 \right) x^{m-n }+ a_{m-n-1 }x^{m-n-1 } + \ldots + a_0 \\
		&= h(x) \in F[x] \\
	.\end{align*} 
	When $\operatorname{deg} h \le m-1<\operatorname{deg} f$, so inductive hypothesis applies, so $h(x) = q_1(x) g(x) + r_1(x)$ where $q_1, r_1 \in F[x]$, $r_1(x) = 0$ or $\operatorname{deg} r_1 < \operatorname{deg} g$.
	But then $f(x)= \left( \frac{a^m}{b^n}x^{m-n}q_1(x) \right)  g(x) + r_1(x) =: q(x)g(x) + r(x)$.
	Where $q(x), r(x) \in F[x]$, and $q, r$ satisfies needed properties. This complete inductive step.
\end{proof}
\begin{remark}
	One can also prove by linear algebra.
\end{remark}
